% SIGIL True Instant Finality — design proposal v0
% Research pass: read crates/sigil-dagknight, sigil-node/producer_signing.rs,
% sigil-braidpool/committee.rs, sigil-net live topic wiring; cross-checked
% against current (2023-2024) BFT/DAG-BFT literature (GRANDPA, Casper FFG,
% HotStuff-2/Jolteon, Mysticeti, SoK: DAG-based Consensus Protocols).
% Build: pdflatex x2.
\documentclass[10pt,a4paper]{article}

\usepackage[margin=2.6cm]{geometry}
\usepackage[T1]{fontenc}
\usepackage{lmodern}
\usepackage{amsmath,amssymb,amsfonts}
\usepackage{xcolor}
\usepackage{booktabs}
\usepackage{enumitem}
\usepackage{microtype}
\usepackage[colorlinks=true,linkcolor=sigilviolet,urlcolor=sigilcyan,citecolor=sigilviolet]{hyperref}
\usepackage{xurl}
\usepackage{longtable}

\definecolor{sigilviolet}{HTML}{8B5CF6}
\definecolor{sigilgold}{HTML}{B7791F}
\definecolor{sigilcyan}{HTML}{0E7490}
\definecolor{ink}{HTML}{1A1428}
\definecolor{slate}{HTML}{D7DCE1}
\definecolor{warn}{HTML}{B91C1C}

\usepackage{tcolorbox}
\tcbuselibrary{skins}
\newcommand{\layman}[3]{%
\begin{tcolorbox}[enhanced,
  interior style={left color=#1!9,right color=white},
  colframe=#1!45!slate,boxrule=0.5pt,arc=3pt,
  borderline west={2.5pt}{0pt}{#1},
  left=9pt,right=9pt,top=7pt,bottom=7pt]
{\scriptsize\textcolor{#1}{\textsc{in plain words}}\; — \;\textbf{\small #2}}\\[3pt]
\small #3
\end{tcolorbox}}
\newcommand{\warnbox}[1]{%
\begin{tcolorbox}[enhanced,
  interior style={left color=warn!8,right color=white},
  colframe=warn!55!slate,boxrule=0.5pt,arc=3pt,
  borderline west={2.5pt}{0pt}{warn},
  left=9pt,right=9pt,top=7pt,bottom=7pt]
{\scriptsize\textcolor{warn}{\textsc{honest risk}}}\\[3pt]
\small #1
\end{tcolorbox}}

\hypersetup{
  pdftitle={SIGIL True Instant Finality — Design Proposal v0},
  pdfauthor={Rocky (Claude Sonnet 5) for SIGIL / sigilgraph},
  pdfsubject={A BFT finality-gadget design layered on DagKnight, grounded in 2023-2024 research},
  pdfkeywords={SIGIL, BFT, finality, GRANDPA, HotStuff-2, Mysticeti, DagKnight, GHOSTDAG}
}

\title{\textcolor{ink}{\LARGE SIGIL True Instant Finality}\\[5pt]
\large A BFT Finality Gadget Layered on DagKnight\\[4pt]
\normalsize Design proposal v0 \textperiodcentered\ 2026-08-25}
\author{Rocky (engineering agent) \textperiodcentered\ for operator review}
\date{}

\begin{document}
\maketitle

\layman{sigilviolet}{What this document is}{%
Right now a SIGIL block only becomes ``safe to trust'' after 512 more blocks are built on top of it
(roughly the chain's confirmation wait) — and even then it is a \emph{probabilistic} guarantee, not a
mathematical one: the code's own comment admits no finite wait is safe against an unbounded adversary.
\textbf{True instant finality} means a block becomes permanently irreversible the instant a quorum of
trusted signers vote for it — zero chance of ever being reversed, by construction, not by odds. This is a
real multi-week engineering project. This document is a concrete design for it, grounded in how
production systems (Ethereum, Polkadot, Sui) actually solve the same problem, and in what this codebase
already has today. \textbf{No code has been written.} This is the design to react to before any is.}

\tableofcontents

\section{The core recommendation}

Split block \emph{production} from block \emph{finality} into two independent layers.

\begin{itemize}[leftmargin=1.4em]
  \item \textbf{Layer 1 — production, unchanged.} DagKnight keeps mining, ordering, and coloring blocks
  exactly as it does today. Nothing about how a block gets made changes.
  \item \textbf{Layer 2 — finality, new.} A small, known, fixed set of signers (\emph{validators})
  periodically look at what DagKnight has already produced and cryptographically \emph{vote}: ``I agree
  this specific point in the order is correct.'' Once enough of them (a \emph{quorum}) sign the
  \emph{same} thing, that point is \textbf{finalized} — locked forever. No future block, no
  longer-looking competing history, nothing, can ever overturn it again.
\end{itemize}

\layman{sigilcyan}{Why two layers instead of one}{%
Today, SIGIL uses one mechanism (DagKnight's block-count wait) to do both jobs at once: decide the
order, \emph{and} decide when to stop worrying about it changing. The two-layer split is how every
modern chain that wants real finality solves this: \textbf{Ethereum} runs Casper FFG (voting) on top of
GHOST (production, DagKnight's closest academic cousin); \textbf{Polkadot} runs GRANDPA (voting) on top
of BABE (production). Polkadot's own docs put it exactly right: ``By using separate protocols for block
production and finality, Polkadot can achieve rapid block creation and strong guarantees of
finality \dots\ GRANDPA guarantees that once a block is finalized, it can never be reverted.''
Production never has to change, slow down, or become more centralized to get real finality — the voting
layer sits on top, watching, not steering.}

\subsection{Why not copy Sui's Mysticeti instead}

Mysticeti (live on Sui mainnet since mid-2024) is the most relevant \emph{recent} DAG-native BFT system,
and deserves the direct comparison the operator asked for. It is \textbf{leaderless and uncertified}:
validators build the DAG by referencing each other's blocks directly; a block is considered decided once
enough later blocks reference it — there are no separate signed vote messages at all. Per the technical
writeup: ``Mysticeti advances the DAG by receiving $n-f$ messages, not by certifying $n-f$ signatures,''
committing in three network round-trips with sub-second latency at real scale on Sui mainnet.

The catch: in Mysticeti, \textbf{the DAG \emph{is} the vote}. Production and finality are not
separable layers — the validator committee \emph{is} the set of block producers, and their act of
referencing each other's blocks \emph{is} the voting. Retrofitting this onto SIGIL would mean replacing
DagKnight's open, PoW-gated production model with a closed, committee-authored DAG — ending permissionless
mining entirely. That is a legitimate, much larger future direction; it is not the ``add finality on top
of what exists'' project requested here. A cross-check against the field's own survey literature (a
2024 SoK on DAG-based consensus) found no published design that hybridizes GHOSTDAG-style probabilistic
fork-choice with a Mysticeti-style uncertified DAG-BFT layer — this proposal's split-layer shape is
closer kin to GRANDPA/Casper~FFG, deliberately.

What \emph{is} borrowed from the newer literature: the \textbf{quorum-certificate mechanics}. HotStuff-2
(2023) showed that a classic three-phase BFT vote can be safely simplified to two phases — ``vote,
then a quorum certificate \emph{is} the commit'' — with no loss of safety, because DagKnight (not the
finality gadget) already handles proposing/ordering. Its production deployment, Jolteon/AptosBFT,
reports sub-second finality under normal conditions using exactly this simplification. This proposal
uses HotStuff-2's two-phase quorum rule, GRANDPA/Casper~FFG's layering shape, and treats Mysticeti's
commit-rule idea (fast implicit commit via later references) as a candidate future optimization, not
the base design.

\section{Validator set}

\subsection{The codebase already answers ``how many validators are enough''}

\texttt{sigil-braidpool::committee} already implements the standard Byzantine-quorum formula:
\begin{align*}
  f_{\max}(n) &= \left\lfloor \frac{n-1}{3} \right\rfloor & \text{(tolerable Byzantine/offline validators)}\\
  \text{quorum}(n) &= n - f_{\max}(n) & \text{(signatures needed to finalize)}\\
  \text{bft\_active}(n) &\iff n \geq 4 & \text{(minimum for BFT to mean anything)}
\end{align*}

\begin{center}
\begin{tabular}{@{}rrrl@{}}
\toprule
$n$ (validators) & $f$ (tolerable) & quorum needed & real BFT? \\
\midrule
1 & 0 & 1 & No — ``trust the one signer,'' not BFT \\
\textbf{2} (Epsilon + happysrv, today) & 0 & 2 (both) & \textbf{No} — and either going offline halts finality \\
\textbf{4} & \textbf{1} & \textbf{3} & \textbf{Yes} — the code's own documented floor \\
7 & 2 & 5 & Yes, more headroom \\
10 & 3 & 7 & Yes, more headroom \\
\bottomrule
\end{tabular}
\end{center}

\warnbox{SIGIL's real, current, live producer count is \textbf{2} (Epsilon and \texttt{happysrv}) — below
the codebase's own minimum for a meaningful Byzantine-fault-tolerance claim. ``Just let today's producers
vote'' is not viable as-is: a 2-of-2 quorum halts finality the instant either node has a routine
restart, and offers no real safety margin against a single compromised key. Reaching $n\!\geq\!4$ needs
at least two more independent machines under operator control.}

\subsection{Three options, weighed against where SIGIL actually is}

\begin{center}
\begin{tabular}{@{}p{2.6cm}p{5.6cm}p{5.7cm}@{}}
\toprule
\textbf{Option} & \textbf{What it means} & \textbf{Fit for SIGIL today} \\
\midrule
A. Fixed, operator-configured & Operator directly picks $n{=}4$–5 validator keys, config-pinned, changed
only by a deliberate coordinated action & \textbf{Recommended.} Simplest; honest about SIGIL's current
operator-run reality; matches the existing \texttt{SIGIL\_TRUSTED\_PRODUCER\_ID\_HEX} ``every node
agrees out-of-band'' precedent already in this codebase \\
B. Stake-weighted / open registration & Anyone stakes or registers a key; voting power proportional &
Not yet — needs a real sybil defense first, and this codebase already has an unresolved, related risk
(fork-choice weighted by raw block count, not real work) that opening validator registration would
extend into the safety layer, not just production \\
C. Delegated / elected committee & Token-holders elect a small validator committee & Reasonable
\emph{future} step once there is a real, widely-held token to elect from — not now \\
\bottomrule
\end{tabular}
\end{center}

\textbf{Recommendation: Option A}, $n{=}4$ or $5$, Ed25519 keys (reusing the exact signing scheme already
live for producer authentication), pinned in a config every node loads at startup.

\section{Message format \& wire protocol}

Reuse Ed25519 for every vote (cheap, already wired in \texttt{producer\_signing.rs}); reuse the
existing periodic hybrid SQIsign5+Ed25519 checkpoint scheme (already built, every 128 blocks) to
additionally post-quantum-attest only the deepest finality anchors, not every vote — SQIsign5 costs
$\sim$1.16s per signature by the code's own measurement, too slow for per-vote use.

\begin{quote}\small\ttfamily
struct FinalityVote \{\\
\quad height: u64,\\
\quad spine\_block\_hash: BlockHash,\ \ \ // DagKnight's own selected spine at that height\\
\quad order\_hash: [u8; 32],\ \ \ \ \ \ \ \ \ // Braid's existing frozen\_acc order commitment\\
\quad validator\_id: ValidatorId,\\
\quad signature: SignatureBytes,\\
\}
\end{quote}

A \emph{certificate} is simply the set of votes agreeing on the same
$(\text{height}, \text{spine\_block\_hash}, \text{order\_hash})$ once that set reaches quorum size.
Propagation reuses the existing gossip-topic convention (\texttt{/sigil/g0/blocks},
\texttt{/sigil/g0/peer-heights}, \ldots) with one new topic, \texttt{/sigil/g0/finality-votes}.

\warnbox{Before any Phase~2 code is written: confirm which gossipsub subscription list the \emph{live}
\texttt{sigil-node} binary actually builds from. \texttt{sigil-net}'s own topic-constant list has a
documented dead-code trap — a same-named constant lives in a different crate (\texttt{flux\_p2p}) and is
the one actually wired in; this exact class of mistake already shipped once (v7.1.44) and quietly did
nothing. Verify before wiring, don't assume.}

\section{Safety and liveness, in plain terms}

\textbf{The one-sentence proof.} If two \emph{different} things both somehow got finalized, each would
need its own quorum of $2f{+}1$ signatures out of $n = 3f{+}1$ total. Two such quorums must overlap in
at least $(2f{+}1)+(2f{+}1)-(3f{+}1) = f{+}1$ validators — meaning at least one \emph{honest} validator
(since at most $f$ are dishonest) would have had to sign both conflicting things. An honest validator
never does that by construction. This is the standard Byzantine quorum-intersection argument underlying
GRANDPA, Casper~FFG, and HotStuff-2 alike — not new math, but the exact reason ``irreversible'' is a real
claim here and not marketing language, \emph{as long as the $\leq f$ assumption holds}.

\begin{center}
\small
\begin{tabular}{@{}p{4.6cm}p{9.4cm}@{}}
\toprule
\textbf{What happens} & \textbf{Consequence} \\
\midrule
$\leq f$ validators go offline / partition & Finality \emph{pauses} for that checkpoint; production keeps
going untouched; catches up once reconnected. A \textbf{liveness} hiccup, never a safety break. \\
$> f$ validators go offline at once & Finality \textbf{halts entirely} until enough return. Nothing
unsafe happens, but nothing new finalizes — a direct, real cost of a small $n$, and not hypothetical
given this fleet's own operational history (Delta down a month; Gamma/Beta permanently gone). \\
A validator equivocates ($\leq f$ dishonest) & Quorum-intersection holds; no conflicting certificate can
form. Both signed votes are visible on the wire — cheap, strong evidence for a future slashing/eviction
mechanism. \\
$> f$ actually dishonest (assumption broken) & \textbf{No BFT design of any kind can prevent this} —
mitigation is entirely about keeping $n$ small enough that the operator directly knows every key
(Option~A) and making the set genuinely hard to expand quietly. \\
\bottomrule
\end{tabular}
\end{center}

\warnbox{This gadget makes whatever DagKnight \emph{already selected} permanent — it does not audit or
correct that selection. If the existing count-based (not difficulty-weighted) blue-score metric can be
gamed to bias which block becomes the spine, the finality gadget will lock in a gamed result exactly as
confidently as an honest one, \emph{forever}. Today a bad selection is only temporary — outrunnable by
enough honest work. After finality, it would not be. Treat fixing blue-score to weigh real work, not raw
block count, as a prerequisite to seriously evaluate before Phase~3, not a nice-to-have.}

\section{Phased implementation plan}

\begin{center}
\small
\begin{tabular}{@{}p{1.6cm}p{4.6cm}p{3.4cm}p{1.6cm}p{2.8cm}@{}}
\toprule
\textbf{Phase} & \textbf{Ships} & \textbf{Consensus effect} & \textbf{Size} & \textbf{Risk if wrong} \\
\midrule
1 & Vote/certificate types, Ed25519 sign/verify, quorum-intersection property tests & None — pure library
code & Days & Low — caught by tests \\
2 & P2P gossip topic, local vote tallying, ``would-be finalized height'' logged alongside today's
estimate; \texttt{Braid::insert()} untouched & None — purely observational, same posture as the existing
frontier-debug diagnostic & 1–2 wks & Low — worst case is bad logs \\
3 & \texttt{Braid} actually gated on real certificates, behind an opt-in flag, isolated test network only
& \textbf{This is the consequential phase} & Weeks & \textbf{Chain-safety bug if wrong} \\
4 & Live activation on sigil-g0, height-gated, after Phase 3 survives adversarial soak testing & Full —
real funds depend on correctness & Smaller if Phase 3 was thorough & Irreversible by design if wrong \\
\bottomrule
\end{tabular}
\end{center}

\section{Honest risk assessment}

A bug in ordinary block production is a rollback-and-fix situation. A bug in a \emph{finality} gadget is
categorically worse, because the entire point of the feature is the promise that nothing can be undone —
getting it wrong either permanently locks in a wrong outcome, or forces breaking the very promise the
feature exists to make, which is worse than never shipping it.

Before Phase~3 ever touches the live chain, this proposal insists on:
\begin{enumerate}[leftmargin=1.4em]
  \item A dedicated multi-node test network, $n{\geq}4$ validators on \emph{genuinely separate}
  machines, run for real wall-clock days, not a CI pass.
  \item Explicit adversarial scenarios: kill a validator mid-vote repeatedly; force one to equivocate
  deliberately; partition the validator set so neither side can reach quorum; replay the exact
  reorder-depth scenario that already forced the $64{\to}512$ \texttt{final\_depth} bump once before.
  \item A genuine go/no-go conversation with the operator before Phase~4 — finality is a bigger trust
  commitment than the existing ``adding a second producer deserves its own explicit go-ahead'' norm
  already practiced in this codebase, and deserves at least the same bar.
  \item The count-vs-work blue-score gap resolved or consciously, explicitly accepted first.
  \item A written-down-in-advance answer for ``what do we do if Phase 3/4 finds a safety bug after
  something is already finalized live'' — decided before activation, not improvised after.
\end{enumerate}

\section*{Decision points needing the operator}

\begin{enumerate}[leftmargin=1.4em]
  \item Go/no-go on Option~A (fixed, small, operator-controlled validator set) versus waiting for a real
  stake/token distribution to justify Option~B/C.
  \item Where do the 2+ additional validator machines come from — new independent boxes/VPS instances
  under operator control.
  \item Fix the blue-score count-vs-work gap before Phase~3, or explicitly accept it and proceed in
  parallel.
  \item Confirm whether producer-signing (\texttt{SIGIL\_PRODUCER\_SIGNING\_SEED\_HEX}) is already
  active on live Epsilon — a config file suggests it might be, contradicting a code comment that says
  it is dormant everywhere.
\end{enumerate}

\end{document}
